% !TeX program = pdflatex
%
% pk-validation-dossier.tex
%
% Build:  pdflatex pk-validation-dossier.tex   (run 3x: ToC, then cross-refs)
%
% Every numeric table in this document is generated by
% .pk_analysis/dossier_tables.py, which drives the shipped constants through the
% same summation the app performs. Regenerate after any constant change.

\documentclass[11pt,a4paper]{article}

\usepackage[T1]{fontenc}
\usepackage[utf8]{inputenc}
\usepackage[margin=2.4cm]{geometry}
\usepackage{amsmath}
\usepackage{amssymb}
\usepackage{booktabs}
\usepackage{longtable}
\usepackage{array}
\usepackage{xcolor}
\usepackage{enumitem}
\usepackage{fancyhdr}
\makeatletter\def\input@path{{../}{../../design/brand/}}\makeatother
\usepackage{isocline}   % docs/research/isocline.sty — Isocline Research template
\usepackage[hidelinks]{hyperref}

% Legacy markup, now neutral: Isocline reports observations, not verdicts.
\newcommand{\bad}[1]{\obs{#1}}
\newcommand{\good}[1]{#1}
\newcommand{\code}[1]{\texttt{\small #1}}

\isoclineshorttitle{PK Model Validation}

\title{\textbf{Isocline Pharmacokinetic Model}\\[0.3em]
       \large Validation Dossier: supported medications, model derivation,\\
       accuracy assessment, and known limitations}
\author{Engineering validation pass}
\date{31 July -- 1 August 2026}

\begin{document}
\maketitle

\begin{abstract}
\noindent
This document describes every medication and route Isocline models, how each
prediction is derived, what published data each constant is anchored to, and how
accurate the resulting predictions are. It is written to be checkable: every
figure quoted is regenerated by a script in \code{.pk\_analysis/}, and the
provenance of every constant is stated, including where that provenance is weak.

The headline results of this validation pass are: two constant-lookup defects that
between them overpredicted by roughly $1{,}090\times$, $5{,}300\times$ and
$344\times$ on three medication/route combinations, plus a third defect in the
chart's forward projection; a testosterone cypionate recalibration against a
primary population-PK source; a conservation-of-mass correction that brought the
three remaining testosterone esters into agreement with it; and a set of
structural limitations that are quantified here rather than fixed, because fixing
them requires a different model rather than different numbers.

Two of those defects were found by running synthetic user scenarios
(Section~\ref{sec:defects}) rather than by inspecting constants, which is the main
methodological conclusion of the exercise.

Sections~\ref{sec:limitations} and~\ref{sec:openquestions} are the ones to read
if you are deciding whether to trust a particular number. Nothing in this document
should be read as clinical guidance.
\end{abstract}

\tableofcontents
\clearpage

% ======================================================================
\section{Scope: what the app models}

Isocline predicts serum hormone concentration over time from a user's logged dose
history. Two medication classes are supported, each across three administration
routes.

\begin{table}[h]
\centering
\small
\begin{tabular}{@{}llll@{}}
\toprule
\textbf{Medication} & \textbf{Route} & \textbf{Ester / form} & \textbf{Output unit} \\
\midrule
Estradiol     & Injectable   & Valerate (EV), Cypionate (EC) & pg/mL \\
Estradiol     & Oral         & (route pseudo-ester)          & pg/mL \\
Estradiol     & Cream        & (route pseudo-ester)          & pg/mL \\
Testosterone  & Injectable   & Cypionate (TC), Enanthate (TE), & ng/dL \\
              &              & Propionate (TP), Undecanoate (TU) & \\
Testosterone  & Oral         & undecanoate (route pseudo-ester) & ng/dL \\
Testosterone  & Cream        & unesterified gel (route pseudo-ester) & ng/dL \\
\bottomrule
\end{tabular}
\caption{The ten medication/route/form combinations the model covers.}
\end{table}

\paragraph{Route pseudo-esters.} The two non-injectable cases in each ester
enumeration describe a \emph{route}, not a chemical ester. They are resolved from
the dose record's administration method by
\code{MedicationManager.effectiveEsterName}, never chosen by the user. That
resolution was introduced in this pass and is the fix for the largest defect
described in Section~\ref{sec:defects}.

\paragraph{What is modeled separately.} The engine returns the \emph{exogenous}
contribution only. Endogenous production is added by the app as a separate
\code{baseLevel} term that decays as gonadal suppression progresses. This matters
for interpreting several constants below, most importantly testosterone
cypionate's half-life.

% ======================================================================
\section{The model}

\subsection{Governing equation}

Each administered dose contributes a Bateman biexponential --- first-order
absorption from a depot into a single well-mixed compartment, first-order
elimination from it:

\begin{equation}
C_i(t) = D_i \cdot \mathrm{EDF} \cdot S \cdot
         \frac{k_a}{k_a - k_e}\left(e^{-k_e t} - e^{-k_a t}\right),
         \qquad t \ge 0
\label{eq:bateman}
\end{equation}

where $D_i$ is the administered dose in mg, $\mathrm{EDF}$ is the effective dose
factor (Section~\ref{sec:edf}), $S$ is a scaling constant with units of
concentration per mg, and $k_a,k_e$ are rate constants in day$^{-1}$. Total
predicted level at time $t$ is the superposition of all doses taken before it,
plus the endogenous baseline, times a user calibration factor:

\begin{equation}
C(t) = \alpha\left[C_{\text{base}}(t) + \sum_{i \,:\, t_i \le t} C_i(t - t_i)\right]
\end{equation}

Superposition is exact here because Equation~\ref{eq:bateman} is linear in dose.
That linearity is also the model's principal simplification: there is no
saturable protein binding, no SHBG dynamics, no dose-dependent clearance.

\paragraph{The $k_a \approx k_e$ limit.} When $|k_a - k_e| < 10^{-4}$
Equation~\ref{eq:bateman} is singular, and the engine substitutes its limit,
$D \cdot \mathrm{EDF} \cdot S \cdot k_a\, t\, e^{-k_a t}$.

\paragraph{Transdermal.} The cream/gel route does not use
Equation~\ref{eq:bateman}. It is modeled as 24 hourly micro-doses of
$D/24$ each, decaying single-exponentially, approximating a 24-hour square-wave
release. Section~\ref{sec:limitations} explains why this under-models intraday
variation.

\subsection{Time to peak}

Differentiating Equation~\ref{eq:bateman} gives

\begin{equation}
t_{\max} = \frac{\ln(k_a/k_e)}{k_a - k_e}
\label{eq:tmax}
\end{equation}

Note that Equation~\ref{eq:tmax} is \emph{symmetric} under $k_a \leftrightarrow
k_e$. This matters more than it looks: for depot esters the two exponentials are
not separately identifiable from clinical data, so which one is labelled
``absorption'' and which ``elimination'' is a modelling convention rather than a
finding. The observable curve is unchanged by the swap, with the amplitude
difference absorbed into $S$.

\subsection{Flip-flop kinetics}

For depot esters, absorption from the oil vehicle is slower than elimination of
the released hormone ($k_a \ll k_e$). The terminal slope of the curve is then
governed by \emph{absorption}, not elimination --- ``flip-flop'' kinetics. The
practical consequence is that the field named \code{halfLifeDays} in the code sets
$k_a$ and controls the entire inter-dose decay, while
\code{eliminationHalfLifeHours} sets $k_e$, which mostly controls where the peak
falls. The names are historical and slightly misleading; the doc comments say so.

\subsection{Effective dose factor}
\label{sec:edf}

$\mathrm{EDF} = \mathrm{EWF} \times F$ converts administered mg of ester into
systemically available mg of active hormone.

The \emph{ester weight factor} is pure stoichiometry --- the mass fraction of the
ester molecule that is the parent hormone:
\[
\mathrm{EWF} = \frac{M_{\text{hormone}}}{M_{\text{ester}}}
\]
with $M_{\mathrm{E}_2} = 272.4$ and $M_{\mathrm{T}} = 288.4$ g/mol. These are the
only constants in the entire model that are not fitted to anything and cannot be
wrong.

$F$ is route bioavailability. Injectable depot $F = 1.0$; oral estradiol
$F = 0.05$ (first-pass metabolism); transdermal estradiol $F = 0.15$; transdermal
testosterone $F = 0.10$, which the AndroGel 1\% label states outright in
\S12.3. Oral testosterone undecanoate carries $F = 1.0$ not because it is fully
bioavailable --- it is not --- but because its absolute bioavailability is not
established in the source. Since $F$ and $S$ are confounded (only their product is
identified by the published data), folding it into $S$ avoids inventing a number.
Predictions are unaffected by where the split is drawn.

\subsection{The scaling constant}

$S$ is an empirical surrogate for $F/V_d$. It is the one constant with no direct
physical interpretation, and historically the one most often wrong. This pass
replaced curve-fitting for $S$ with closed forms wherever the sources allowed;
see Sections~\ref{sec:tc} and~\ref{sec:identity}.

% ======================================================================
\section{Encoded constants}

\begin{table}[h]
\centering
\small
\begin{tabular}{@{}lrrrrrrr@{}}
\toprule
& \multicolumn{2}{c}{\textbf{slow exponential}} & \textbf{fast} & & & & \\
\cmidrule(lr){2-3}
\textbf{Form} & $t_{1/2}$ (d) & $k_a$ (d$^{-1}$) & $k_e$ (d$^{-1}$)
              & EWF & $F$ & $S$ & $t_{\max}$ (d) \\
\midrule
EV            & 2.6000 & 0.2666 & 0.7737 & 0.764 & 1.00 & 396.40 & 2.10 \\
EC            & 6.4700 & 0.1071 & 0.4310 & 0.687 & 1.00 & 273.00 & 4.30 \\
E2 oral       & 0.0791 & 8.7629 & 1.3859 & 1.000 & 0.05 & 1250.00 & 0.25 \\
E2 cream      & 1.5000 & 0.4621 & 5.5452 & 1.000 & 0.15 & 155.00 & --- \\
\midrule
TC            & 3.8390 & 0.1806 & 1.2196 & 0.699 & 1.00 & 67.13 & 1.84 \\
TE            & 7.0000 & 0.0990 & 8.3178 & 0.720 & 1.00 & 457.67 & 0.54 \\
TP            & 1.0000 & 0.6931 & 8.3178 & 0.837 & 1.00 & 457.67 & 0.33 \\
TU            & 20.0000 & 0.0347 & 8.3178 & 0.631 & 1.00 & 457.67 & 0.66 \\
T oral (TU)   & 0.0900 & 7.7016 & 8.3053 & 0.631 & 1.00 & 13.57 & 0.13 \\
T gel         & 0.5784 & 1.1984 & ---    & 1.000 & 0.10 & 132.98 & --- \\
\bottomrule
\end{tabular}
\caption{Constants as shipped, and the $t_{\max}$ each implies via
Equation~\ref{eq:tmax}. The cream routes have no $t_{\max}$ because they are not
modeled with Equation~\ref{eq:bateman}; E2 cream's $k_e$ is stored but unused.}
\label{tab:constants}
\end{table}

\paragraph{Reading Table~\ref{tab:constants}.} Three things in it are worth
noticing before the derivations.

TC's slow half-life (3.84 d) is now \emph{shorter} than TE's (7.0 d), inverting
the previous ordering even though the two esters are clinically near
interchangeable. This is an artifact of TC being the only ester with a
primary-source refit, not a claim about relative clearance.

TE, TP and TU share an identical $S$. That is not a copy-paste error; it follows
from the identity in Section~\ref{sec:identity} and the fact that all three still
share a $k_e$.

E2 oral's stored half-lives were first \emph{restated} (0.5\,d $\to$ 0.2773\,d,
3.0\,h $\to$ 12.0035\,h) without changing any prediction. The engine had never
read those fields for oral doses; it hardcoded $k_a = 2.5$ and $k_e = 1.386$
inline, and $S = 1250$ was fit against those. Restating them made the constants
stop disagreeing with the code --- and, once they were visible, made it possible
to see that $k_a$ was wrong.

It has since been \emph{refit}: 0.2773\,d $\to$ 0.0791\,d, i.e.\ $k_a = 2.5 \to
8.762$\,d$^{-1}$, moving $t_{\max}$ from 12.7\,h into the published window at
6.0\,h. Unlike the restatement, this one moves predictions. $k_e$ and $S$ are
unchanged; the reasoning for holding $S$ is in Section~\ref{sec:limitations}.

% ======================================================================
\section{Derivations by form}

\subsection{Estradiol valerate and cypionate}

Recalibrated on 29 July 2026 (before this pass) against
transfemscience.org's compiled tables, which aggregate published single-dose and
steady-state estradiol PK. $k_e$ was fit so Equation~\ref{eq:tmax} reproduces the
published time-to-peak; $k_a$ was fit so the modeled steady-state trough
reproduces Table 10's 5\,mg-weekly figures; $S$ was anchored to the same table's
steady-state peak.

The important decision recorded there, and re-made in this pass for testosterone,
is that single-dose and steady-state targets were \textbf{not simultaneously
satisfiable}, and steady state was chosen. Section~\ref{sec:openquestions}
returns to this.

\subsection{Testosterone cypionate}
\label{sec:tc}

TC is the only form in the model derived from a primary population-PK
publication: Bi et al.\ (2018), \emph{CPT: Pharmacometrics \& Systems
Pharmacology}, \texttt{doi:10.1002/psp4.12287}. Table 2 of that paper reports
$\mathrm{CL}/F = 2.6$\,kL/day, $V/F = 14.4$\,kL, and $K_a = 1.22$\,day$^{-1}$.

\paragraph{Rate constants.} The disposition rate is
$\mathrm{CL}/V = 2.6/14.4 = 0.1806$\,day$^{-1}$, giving a half-life of 3.839\,d,
which the code stores as \code{halfLifeDays}. $K_a = 1.22$ gives
$\ln 2 / 1.22 = 13.64$\,h, stored as \code{eliminationHalfLifeHours}. (Per the
symmetry noted after Equation~\ref{eq:tmax}, the labels are conventional.)

\paragraph{Why 3.84 days and not 7--8.} The commonly quoted TC half-life is
around a week, and Bi addresses the discrepancy directly. Noncompartmental
analysis of their own data yields 6.9\,d, and their model yields 6.87\,d
\emph{when endogenous secretion is assumed to be zero} --- but 4.05\,d once
endogenous production is accounted for. Because Isocline models endogenous production
separately as \code{baseLevel}, the exogenous-only figure is the one that belongs
in the ester constant. Bi's own median post-hoc estimate of 4.05\,d agrees with
our 3.839\,d to 5.2\%.

\paragraph{Scaling constant, in closed form.} Matching Equation~\ref{eq:bateman}
to the one-compartment model underlying Bi's parameters gives $S$ analytically
rather than by curve fit:
\begin{equation}
S = \frac{10^5 \cdot K_a}{V_L \cdot \mathrm{EDF} \cdot k_e}
  = \frac{10^5 \cdot 1.22}{14400 \times 0.699 \times 0.180556} = 67.13
\label{eq:tcS}
\end{equation}
The $10^5$ converts mg/L to ng/dL; EDF appears because the app multiplies the dose
by it while Bi's $V/F$ already absorbs ester conversion. This reproduces Bi's
population model exactly, which is why the corresponding row of
Table~\ref{tab:accuracy} shows 0.0\% error --- that row is a consistency check,
not evidence.

\paragraph{Resolution of a standing concern.} The project's own
\code{OPEN\_CONCERNS.md} \S2 had flagged the encoded TC $t_{\max}$ window
(1--2\,d) as probably wrong, because a secondary source placed $C_{\max}$ on day
4--5. The primary source settles it in favour of the encoded window: Bi's
parameters give $t_{\max} = 1.84$\,d via Equation~\ref{eq:tmax}, and propagating
the published 95\% CI on $K_a$ (0.785--2.67) gives a band of 1.08--2.43\,d, which
never reaches day 4--5. Read TC $t_{\max}$ as $1.84 \pm 0.7$\,d; Bi's own
limitations section (1--2 absorption-phase samples per subject, $K_a$
RSE 44.9\%) explains the imprecision.

\subsection{Testosterone enanthate, propionate, undecanoate}
\label{sec:identity}

None of these three has a primary source in hand. Their scaling constants are
nonetheless \emph{not} free parameters, because of a conservation argument that
follows from the TC refit.

Integrating Equation~\ref{eq:bateman} over all time gives exactly
\[
\int_0^{\infty} \frac{k_a}{k_a-k_e}\left(e^{-k_e t}-e^{-k_a t}\right)\,dt
= \frac{k_a}{k_a-k_e}\left(\frac{1}{k_e}-\frac{1}{k_a}\right) = \frac{1}{k_e}
\]
so at steady state with dosing interval $\tau$, superposition gives
\begin{equation}
C_{\mathrm{avg}} = \frac{D \cdot \mathrm{EDF} \cdot S}{\tau \, k_e}
\label{eq:cavgmodel}
\end{equation}
The true steady-state average is (mg of testosterone per day) divided by
clearance --- and clearance is a property of \emph{testosterone}, not of the ester
carrying it. It is therefore already pinned by the TC refit:
$\mathrm{CL}_T = 2600 \times \mathrm{EWF}_{TC} = 1817.4$\,L/day. Equating the two
expressions, the ester weight factor cancels:
\begin{equation}
\boxed{\;S = \frac{10^5 \cdot k_e}{\mathrm{CL}_T \cdot F}\;}
\label{eq:identity}
\end{equation}

$S$ depends only on $k_e$ and clearance. The ester's half-life --- $k_a$ --- does
not appear at all, which is worth stating plainly because two now-deleted unit
tests asserted an ordering of $S$ by half-life and were reasoning from the wrong
quantity.

\paragraph{The load-bearing check.} TC satisfies Equation~\ref{eq:identity} to
0.03\% \emph{without having been fit to it} --- its $S$ came from
Equation~\ref{eq:tcS}, via $V/F$ and $K_a$. Two independent derivations landing
on the same number is what makes it defensible to apply the identity to the other
three.

\begin{table}[h]
\centering
\small
\begin{tabular}{@{}lrrrlrrr@{}}
\toprule
\textbf{Ester} & $S$ before & $S$ now & $S$ implied & error before
               & \textbf{Regimen} & $C_{\mathrm{avg}}$ before & after \\
\midrule
Cypionate   & 1081.0 & 67.13 & 67.11 & $16.11\times$ & 100\,mg q7d & 1298 & 550 \\
Enanthate   & 923.0 & 457.67 & 457.67 & $2.02\times$ & 100\,mg q7d & 1141 & 566 \\
Propionate  & 359.0 & 457.67 & 457.67 & $0.78\times$ & 50\,mg q2d & 903 & 1151 \\
Undecanoate & 214.0 & 457.67 & 457.67 & $0.47\times$ & 1000\,mg q70d & 232 & 496 \\
\bottomrule
\end{tabular}
\caption{The clearance identity applied. $C_{\mathrm{avg}}$ in ng/dL, exogenous
only. The ``before'' column for cypionate also reflects its $k_e$ moving from
2.0\,h to 13.64\,h.}
\label{tab:identity}
\end{table}

\paragraph{What was wrong.} The previous values spanned $0.47\times$ to
$2.02\times$ of the implied value and were \emph{mutually inconsistent} --- the
app told users on different esters different things about the same clearance.
Enanthate is the clearest case: a standard 100\,mg weekly regimen predicted
$C_{\max} = 1753$\,ng/dL including baseline, a level that reads as clearly
supraphysiological, while contradicting its primary-sourced sibling by a factor
of two.

\paragraph{What this does not fix.} Magnitude only. TE, TP and TU retain a
placeholder 2\,h $k_e$, so their $t_{\max}$ (Table~\ref{tab:tmax}) is badly wrong.
See Section~\ref{sec:limitations}.

\subsection{Transdermal testosterone}

Anchored to the AndroGel 1\% prescribing information, Table 5 (Day 180 steady
state, $n=44$): 50\,mg/day $\to C_{\mathrm{avg}} = 555$\,ng/dL. The terminal
half-life comes from the label's washout statement --- levels return to baseline
in 48--72\,h, implying $\sim$5\% residual at 60\,h, giving $t_{1/2} = 0.5784$\,d.

Two limitations were documented rather than fitted around; both are in
Section~\ref{sec:limitations}.

\subsection{Oral testosterone undecanoate}

Anchored to Swerdloff et al.\ (2020) / the Jatenzo label: 237\,mg twice daily
$\to C_{\mathrm{avg}} = 488.7$\,ng/dL, $t_{\max}$ 2--4\,h. Both rate constants
were solved jointly so that $t_{\max}$ lands at 3\,h and the steady-state
$C_{\max}/C_{\mathrm{avg}}$ comes as close as the model permits to the reported
2.06. Section~\ref{sec:limitations} explains why it cannot reach it.

% ======================================================================
\section{Accuracy}

\subsection{Against the anchors}

\begin{table}[h]
\centering
\small
\begin{tabular}{@{}lllrrl@{}}
\toprule
\textbf{Regimen} & \textbf{Quantity} & \textbf{Model} & \textbf{Published}
                 & \textbf{Error} & \textbf{Source} \\
\midrule
EV 5\,mg q7d & $C_{\max}$ & 384 & 384 & \good{$+0.1\%$} & TFS Table 10 \\
EV 5\,mg q7d & $C_{\min}$ & 142 & 142 & \good{$+0.1\%$} & TFS Table 10 \\
EC 5\,mg q7d & $C_{\max}$ & 339 & 339 & \good{$-0.1\%$} & TFS Table 10 \\
EC 5\,mg q7d & $C_{\min}$ & 262 & 262 & \good{$-0.1\%$} & TFS Table 10 \\
TC 100\,mg q7d & $C_{\mathrm{avg}}$ & 550 & 549 & \good{$+0.0\%$} & Bi Table 2 \\
T gel 50\,mg/d & $C_{\mathrm{avg}}$ & 555 & 555 & \good{$+0.1\%$} & AndroGel Table 5 \\
T gel 75\,mg/d & $C_{\mathrm{avg}}$ & 833 & 601 & \bad{$+38.6\%$} & AndroGel Table 5 \\
T gel 100\,mg/d & $C_{\mathrm{avg}}$ & 1111 & 713 & \bad{$+55.8\%$} & AndroGel Table 5 \\
Oral TU 237\,mg BID & $C_{\mathrm{avg}}$ & 489 & 489 & \good{$-0.0\%$} & Jatenzo \\
Oral TU 237\,mg BID & $C_{\max}$ & 789 & 1008 & \bad{$-21.7\%$} & Jatenzo \\
\bottomrule
\end{tabular}
\caption{Predicted versus published. Rows with near-zero error are the anchors
themselves and demonstrate only that the fit converged. The three coloured rows
are the informative ones; both are explained in Section~\ref{sec:limitations}.
Estradiol in pg/mL, testosterone in ng/dL.}
\label{tab:accuracy}
\end{table}

\subsection{Out-of-sample checks}

Rows in Table~\ref{tab:accuracy} with 0.0\% error prove nothing --- a fit agrees
with its own target by construction. The following comparisons were \emph{not}
used in any fit.

\begin{table}[h]
\centering
\footnotesize
\begin{tabular}{@{}lrrrl@{}}
\toprule
\textbf{Check} & \textbf{Model} & \textbf{Published} & \textbf{Error} & \textbf{Source} \\
\midrule
TC single 200\,mg $C_{\max}$ (exogenous) & 997 & 872 & $+14.3\%$
  & 200\,mg IM lit.\ 1112 less 240 basal \\
TC 100\,mg q7d $C_{\max}$ + basal & 976 & 1150 & $-15.1\%$
  & Bi Figure 2, 100\,mg panel \\
EV single 2\,mg $C_{\max}$ & 119 & 118 & \good{$+1.0\%$} & TFS Table 9 \\
EC single 2\,mg $C_{\max}$ & 59 & 62 & \good{$-5.1\%$} & TFS Table 9 \\
\bottomrule
\end{tabular}
\caption{Independent validation. The estradiol rows are the strongest evidence in
this document: constants fit purely to steady-state data reproduce single-dose
peaks within 1\% and 5\%.}
\end{table}

The two testosterone rows agree to about 15\% in opposite directions, which is
roughly the precision the source itself supports given Bi's reported RSEs. They
are consistency checks, not tight validation.

\subsection{Time to peak}

\begin{table}[h]
\centering
\small
\begin{tabular}{@{}llll@{}}
\toprule
\textbf{Form} & \textbf{Model $t_{\max}$} & \textbf{Published} & \textbf{Source quality} \\
\midrule
EV      & 2.10\,d & 2.1\,d & compiled tables \\
EC      & 4.30\,d & 3.1--4.3\,d & compiled tables \\
E2 oral & \good{6.0\,h} & \good{4--8\,h} & secondary \\
TC      & \good{1.84\,d} & \good{1.84\,d} & \textbf{primary} (Bi Table 2) \\
TE      & \bad{12.9\,h} & \bad{1.5--7\,d} & secondary \\
TP      & \bad{7.8\,h} & \bad{$\sim$0.6\,d} & sparse \\
TU      & \bad{15.9\,h} & \bad{$\sim$7\,d} & secondary (Wang 2010) \\
T oral  & \good{3.0\,h} & \good{2--4\,h} & label \\
\bottomrule
\end{tabular}
\caption{Where the model puts the peak. This is the table that shows what is still
uncalibrated: magnitude is right everywhere, shape is not.}
\label{tab:tmax}
\end{table}

% ======================================================================
\section{Persona validation}

Section-level agreement with published tables is necessary but not sufficient:
those tables all describe clean steady state on a fixed regimen, and real users
miss doses, change dose, change route, and dose on intervals nobody studied.
Eight synthetic users were run end to end through the same summation the app
performs. Full output: \code{python3 .pk\_analysis/personas.py}.

\begin{longtable}{@{}p{2.6cm}p{4.5cm}p{7.4cm}@{}}
\toprule
\textbf{Persona} & \textbf{Regimen} & \textbf{What it tests, and the result} \\
\midrule
\endhead
Maya, 24 & 4\,mg EV IM q7d, 18 months, perfect adherence
 & Control. $C_{\max}$ 307, $C_{\min}$ 114\,pg/mL --- exactly $4/5$ of the 5\,mg
   source figures, as linearity requires. \good{Pass.} \\[0.4em]
Devon, 31 & 100\,mg TC q7d $\to$ 200\,mg q14d at month 6
 & Superposition across a regimen change. $C_{\mathrm{avg}}$ preserved at
   790\,ng/dL as it must be; trough falls $561 \to 381$ while peak rises
   $976 \to 1339$, and peak/trough widens $1.74\times \to 3.51\times$. The model
   reproduces the clinical trade-off and correctly refuses to reward the longer
   interval. \good{Pass.} \\[0.4em]
Priya, 47 & 2\,mg oral E2 BID for 4 months, then 4\,mg EV q7d
 & Route switch; historical doses must keep their own kinetics. They do --- the
   oral segment retains $F = 0.05$ and the fast rates. This persona also found the
   oral $t_{\max}$ defect, which collapsed her modeled intraday swing to
   $1.1\times$; after the $k_a$ refit it is $1.37\times$ at an unchanged
   $C_{\mathrm{avg}}$ of 180\,pg/mL. \good{Pass.} \\[0.4em]
Sam, 19 & 50\,mg/day T gel, misses $\sim$2 days/week
 & The configuration that used to overpredict by three orders of magnitude, plus
   skipped-dose handling. Now $C_{\mathrm{avg}}$ 637\,ng/dL; his missed days cost
   20\% of his average and 63\% of his trough. \good{Pass.} \\[0.4em]
Rowan, 38 & 4\,mg/day compounded E2 cream, onboarded in-app
 & Found a defect: onboarding stored the ester as \code{"Cream"}, matching no
   enum case. \good{Fixed at both ends.} \\[0.4em]
Kai, 52 & 237\,mg oral TU BID, 10-day interruption
 & Washout and re-accumulation on the fastest route. Falls to the endogenous
   floor within a day; no phantom drug survives the gap; recovers to 100\% of
   pre-gap within 2 days. \good{Pass.} \\[0.4em]
Jules, 28 & 100\,mg TE IM q7d
 & Found a defect: steady-state $C_{\mathrm{avg}}$ was $2\times$ too high and
   contradicted primary-sourced TC. \good{Fixed} by Equation~\ref{eq:identity}
   ($C_{\max}$ 990, $C_{\min}$ 637\,ng/dL). Shape still wrong --- $t_{\max}$
   0.46\,d. \\[0.4em]
Nadia, 35 & 5\,mg EC q10d --- an interval nobody studied
 & Extrapolation outside the calibration envelope. Trough holds in band but sits
   40\% below the same dose at q7d. Structurally sound but inherits every
   assumption of the fit. \good{Pass, with caveat.} \\
\bottomrule
\end{longtable}

All eight land inside their clinical reference bands. Two of the eight found
defects that inspection had missed, which is the argument for running scenarios
rather than only checking constants.

% ======================================================================
\section{Defects found and fixed}
\label{sec:defects}

\subsection{Non-injectable routes fell through to hardcoded constants}

The scaling-constant lookup inside \code{estimatedLevels} was guarded on
\code{medication == "Estrogen"}. A \emph{testosterone} dose on the Oral or Cream
route therefore matched no branch and fell through to the hardcoded literals
\code{12000.0} and \code{8000.0} --- which are the pre-28-July \emph{estrogen}
values that the codebase already documents as $9.6\times$ and $52\times$ too
large even for estrogen.

Three errors compounded: the wrong $S$; the wrong EDF (resolved from the stored
injectable ester rather than the route, so a gel dose used cypionate's
$0.699 \times 1.0$ instead of the gel's $1.0 \times 0.10$); and a hardcoded
absorption rate that also changed day-to-day accumulation.

\begin{table}[h]
\centering
\small
\begin{tabular}{@{}lrrrr@{}}
\toprule
\textbf{Scenario} & \textbf{Before} & \textbf{After} & \textbf{Factor} & \textbf{Published} \\
\midrule
T gel 50\,mg/d & 605{,}183 & 555 & $1{,}090\times$ & 555 \\
Oral TU 237\,mg BID & 2{,}589{,}554 & 489 & $5{,}299\times$ & 489 \\
E2 cream 4\,mg/d, onboarded & 69{,}263 & 201 & $344\times$ & --- \\
\bottomrule
\end{tabular}
\caption{Steady-state $C_{\mathrm{avg}}$, ng/dL for testosterone and pg/mL for
estradiol. These are simulated through both code paths, not estimated from
constant ratios.}
\end{table}

\paragraph{Reachability.} This was not theoretical. The settings screen offered
every administration method for every medication type, while the ester list
returned an empty array for testosterone + Oral\slash\allowbreak Cream --- so the
previously selected injectable ester stayed silently attached to the record and
the user saw no ester picker at all.

\paragraph{Fix.} \code{MedicationManager.effectiveEsterName(ester:route:)} makes
the \emph{route} decide which constants apply, and both \code{pkParams} and
\code{dosePotencyFactor} now go through it. The hardcoded literals and rates are
gone; every route reads its constants from the ester enumeration. The settings
and onboarding ester lists were narrowed to \code{injectableCases}.

\subsection{Onboarding stored an ester that matched nothing}

A cream dose created through onboarding was saved with ester \code{"Cream"} ---
the raw administration-method name, which matches no case in
\code{MedicationEster}. The same lookup failure applied, so estradiol cream users
who onboarded in-app got the hardcoded 8000 while users who configured the
identical route from Settings stored \code{"Transdermal"} and were fine.

Fixed at both ends: onboarding now stores the route pseudo-ester, and the route
normalization above repairs records that already exist.

\subsection{Forward projection assumed a 7-day interval for every injectable}

With exactly one logged dose the chart projects a future cadence so the user can
preview a trajectory. That interval came from the administration method, which
returns 7 days for anything injectable. For TU --- dosed every 10--14 weeks --- a
single logged 1000\,mg dose was projected as \emph{weekly} 1000\,mg injections,
reaching $\sim$2850\,ng/dL against a single-dose $\sim$1190. TP, dosed every other
day, was under-projected for the mirror-image reason.

Fixed by adding \code{typicalIntervalDays} to both ester enumerations. A
user-supplied frequency still takes precedence.

% ======================================================================
\section{Known limitations}
\label{sec:limitations}

These are quantified rather than fixed. Each needs either a source we do not have
or a structurally different model.

\subsection{TE, TP and TU have the right magnitude and the wrong shape}

All three retain a placeholder 2\,h $k_e$, which pins their $t_{\max}$ to
0.3--0.7\,d against published windows of 1.5--7\,d (Table~\ref{tab:tmax}). The
peak arrives days early and the curvature between peak and trough is wrong.
Steady-state average is correct.

They were not fitted because the usable TE steady-state profile is
\emph{subcutaneous} --- which absorbs more slowly and flatter than the
intramuscular route modeled here, so its $k_a$ does not transfer without
justification --- and because that figure comes from a secondary summary rather
than a primary table. TP's published data is sparse.

A visible consequence is deliberately left in place: TU's modeled $t_{\max}$
(0.66\,d) is now \emph{earlier} than TC's (1.84\,d), inverting the clinical
ordering. A unit test asserts this inversion and instructs whoever fixes TU's
$k_a$ to delete the test and restore the clinical assertion. That test passed
before this pass only because TC was also wrong, in the same direction, by more.

\subsection{Testosterone undecanoate needs a three-parameter model}

Over a 70-day interval the two-parameter Bateman produces a peak/trough ratio of
about $11\times$. The real drug is $2$--$3\times$. No choice of $(k_a, k_e, S)$
reproduces that flatness.

The magnitude correction in Section~\ref{sec:identity} improves TU's
$C_{\mathrm{avg}}$ (232 $\to$ 496\,ng/dL exogenous, against a published
$\sim$500--600 total) but makes its $C_{\max}$ overshoot, because correcting
magnitude cannot correct shape. That trade was taken knowingly: $C_{\mathrm{avg}}$
is what clinicians titrate on and what a user lives at, and the peak was already
wrong.

\subsection{Transdermal intraday ripple is structurally under-modeled}

The 24-hourly-micro-dose construction has no absorption phase and tops out at
$C_{\max}/C_{\min} \approx 1.05$ for testosterone gel, against an observed
$2.24$. Reproducing $2.24$ would require a terminal half-life near 1\,h, which
contradicts the label's own 48--72\,h washout by an order of magnitude. The two
cannot be satisfied together in this model. $C_{\mathrm{avg}}$ is exact.

\subsection{Transdermal dose-response is anchored to one arm}

The AndroGel 1\% label's dose-response is markedly sublinear: 50\,mg $\to$ 555,
75\,mg $\to$ 601, 100\,mg $\to$ 713\,ng/dL. The 1.62\% label inverts outright ---
81\,mg $C_{\mathrm{avg}}$ 537 is \emph{below} 60.75\,mg $C_{\mathrm{avg}}$ 643.

This is a titration artifact: higher-dose arms are enriched for poor absorbers,
because subjects were escalated precisely when they responded badly. It is
selection confounding, not saturable pharmacokinetics, and a linear model should
not be bent to reproduce it. $S$ is anchored to the 50\,mg arm --- the standard
starting dose and the least confounded --- so predictions at 75 and 100\,mg run
high by 39\% and 56\% (Table~\ref{tab:accuracy}).

Gel absorption also varies enormously between individuals. The app's calibration
factor against the user's own blood tests is the intended remedy for this and is
more important on this route than on any other.

\subsection{Oral testosterone under-predicts peak fluctuation}

With $t_{\max}$ pinned to the reported 2--4\,h window, the model's achievable
$C_{\max}/C_{\mathrm{avg}}$ tops out at 1.61 against a reported 2.06, so
$C_{\max}$ reads 789 rather than 1008\,ng/dL.

The reported ratio is itself inflated by mean-of-maxima bias: the published
$C_{\max}$ is the mean of per-subject maxima, and $\mathbb{E}[\max] \ge
\max[\mathbb{E}]$, with a reported SD of $\pm 581$ on a mean of 1008. Anchoring
to $C_{\mathrm{avg}}$ --- the label's own primary endpoint --- and accepting a
documented under-prediction of the peak was judged the better trade.

\subsection{Oral estradiol peaked too late --- refit}

$k_a = 2.5$ and $k_e = 1.386$ gave $t_{\max} = 12.7$\,h against a published oral
estradiol $t_{\max}$ on the order of 4--8\,h. On a twice-daily schedule this
mattered more than the number suggests: $t_{\max}$ exceeded the 12\,h dosing
interval, so each dose was still rising when the next arrived and the modeled
curve flattened to a peak/trough of $\sim$1.1$\times$ across the day. Real oral
estradiol swings considerably more.

$k_a$ is now $8.762$\,/day, placing $t_{\max}$ at $6.0$\,h and widening the modeled
BID swing to $\sim$1.37$\times$. Two things in that are choices rather than fits,
and are recorded as such.

First, \emph{where} inside the window. Nothing picks a point in 4--8\,h --- it is
a range from secondary sources, not a distribution with a mode --- so $t_{\max}$
is pinned to the midpoint. Note also that Equation~\ref{eq:tmax} is symmetric
under $k_a \leftrightarrow k_e$, so inverting it for $k_a$ has two roots; the fit
is bracketed on the $k_a > k_e$ branch.

Second, \emph{what to preserve}. Integrating the Bateman term gives
$\mathrm{AUC}/\text{dose} = D^{\mathrm{eff}} S / k_e$, with no $k_a$ in it: $S$
and $k_a$ are coupled through the peak but not through the average. One
$(S, k_a)$ pair therefore cannot hold both the published $t_{\max}$ and the
published $C_{\max}$, which makes ``refit $k_a$ and re-derive $S$'' a hidden
choice of anchor rather than a mechanical step. We anchored on
$C_{\mathrm{avg}}$: 2\,mg daily stays at 90\,pg/mL, unchanged, and it is the
quantity clinicians titrate on. The accepted cost is single-dose $C_{\max}$
rising from 60 to 88\,pg/mL for 2\,mg. Holding $C_{\max}$ at 60 instead would
have required $S = 848.5$ and dropped $C_{\mathrm{avg}}$ to 61\,pg/mL --- moving
every average and trough figure in the app to buy back one single-dose peak. This
is the same trade, and the same direction, as the oral TU decision above and the
injectable re-anchor.

\subsection{Model-wide simplifications}

Dose-linearity with no saturable binding or SHBG dynamics; a single well-mixed
compartment; constant clearance; no covariates (body composition, age, hepatic
function, concomitant medications beyond the one spironolactone adjustment); and
no representation of between-subject variability, which for the transdermal route
in particular is the dominant source of error. Predictions are population means,
not individual forecasts. The per-user calibration factor exists to close some of
that gap once blood tests are available.

% ======================================================================
\section{Open questions requiring non-engineering judgement}
\label{sec:openquestions}

Three things in this model are decisions rather than derivations, and were made by
engineering because they blocked the work. They deserve explicit ratification.

\subsection{Single-dose targets versus steady-state targets}

For both estradiol esters and all four testosterone esters, the single-dose peak
targets written into the code's own documentation and the published steady-state
data were \textbf{mutually unsatisfiable}. For estradiol the gap was large:
holding the app's documented single-dose targets put steady state $+110\%$ (EV)
and $+206\%$ (EC) above the published table.

Steady state was chosen, both times, on the reasoning that it is what users
actually live at and what clinicians titrate on. The consequence is that the
single-dose figures in the doc comments are now \emph{derived} values, not
enforced targets. If the product intent was the reverse, this needs revisiting.

\subsection{Provenance of the original targets}

The estradiol single-dose targets claimed provenance from Düsterberg et al.\ and
Tiong et al., but do not match those sources ($2.1\times$ for EV, $3.1\times$ for
EC). Whoever set them either used a different reference population or assay, or
made an error. That question is unresolved.

Relatedly: the previously shipped EC constant ($S = 2700$) produced 51\,pg/mL for
a 2\,mg dose, not the 180 its own comment claimed --- a $-71.5\%$ mismatch. Every
other ester matched its stated target to within 0.1\%, which suggests that one
value was never checked against the target written beside it. Worth understanding
how that survived review, since the same process produced the others.

\subsection{What the TU target range is describing}

The TU target range written in the code (400--700\,ng/dL) sits well below the
$\sim$1100\,ng/dL $C_{\max}$ reported for a single 1000\,mg injection. It may be
describing an intended steady-state \emph{trough} rather than a single-dose peak.
No numeric change was made pending confirmation of which quantity it means.

\subsection{Communicating the change}

Predictions move substantially for most configurations in this pass --- see
\code{OPEN\_CONCERNS.md} \S5 for the per-form table. Nothing about anyone's actual
hormone levels has changed; the previous numbers were wrong.

No user-facing release note was written, because the app is not yet deployed and
there are no existing users to inform. That decision should be revisited before
launch. If any of these numbers ship and then move again, the tone will matter:
users of this app are tracking a medical transition, and an unexplained apparent
halving of their levels is alarming in a way that a normal changelog entry does
not address.

% ======================================================================
\section{Verification and reproducibility}

\subsection{Guards in the codebase}

The Swift test suite carries the calibration claims as executable assertions.
The load-bearing ones added or rewritten in this pass:

\begin{itemize}[nosep,leftmargin=1.4em]
\item \code{testosteroneScalingSatisfiesClearanceIdentity} --- checks
      Equation~\ref{eq:identity} for every injectable T ester. Fails on drift in
      any of $S$, $k_e$, $F$ or the clearance constant. This is the assertion that
      would have caught the TE defect.
\item \code{cypionateScalingMatchesAnalyticForm} --- checks
      Equation~\ref{eq:tcS} independently, which is what makes the identity
      trustworthy.
\item \code{cypionateTPeakTime} --- pins TC $t_{\max}$ to Bi's 1.84\,d and to the
      CI band implied by $K_a$'s published interval.
\item \code{undecanoatePeakTimeIsUncalibrated} --- asserts the \emph{known-wrong}
      TU/TC ordering, so the gap stays visible instead of being quietly
      re-broken.
\item \code{testosteroneEliminationHalfLife} --- asserts that TE, TP and TU still
      carry the 2\,h placeholder, so nobody assumes they are calibrated.
\end{itemize}

\subsection{Scripts}

\begin{table}[h]
\centering
\small
\begin{tabular}{@{}ll@{}}
\toprule
\textbf{Script} & \textbf{Purpose} \\
\midrule
\code{verify\_constants.py} & Independent check of the \emph{shipped} constants. Run first. \\
\code{personas.py}          & The eight end-to-end user scenarios. \\
\code{fit\_testosterone.py} & Derivation of TC, gel and oral T constants. \\
\code{fit\_constants.py}    & Derivation of the estradiol ester constants. \\
\code{dossier\_tables.py}   & Regenerates every numeric table in this document. \\
\code{pk\_engine\_replica.py} & Python mirror of the Swift engine; all of the above import it. \\
\bottomrule
\end{tabular}
\end{table}

\code{verify\_constants.py} is deliberately \emph{not} the fitting code. It takes
the shipped numbers at face value, drives them through the same summation the app
performs, and compares against the source figures. A fit that agrees with itself
proves nothing; this closes the loop. It also contains a replica of the previous
code path so the defect magnitudes in Section~\ref{sec:defects} are measured
rather than asserted.

\paragraph{Keeping the replica in sync.} If you change a constant in
\code{PharmacokineticConstants.swift}, change it in \code{pk\_engine\_replica.py}
too. Every validation script reads from the replica, so a divergence silently
invalidates all of them.

% ======================================================================
\section{Sources}

\begingroup\sloppy
\begin{description}[leftmargin=0pt,labelsep=0pt,style=nextline,itemsep=0.3em]
\item[Bi Y.\ et al.\ (2018)] ``Population Pharmacokinetic/Pharmacodynamic
  Modeling of Depot Testosterone Cypionate in Healthy Male Subjects.''
  \emph{CPT: Pharmacometrics \& Systems Pharmacology}.
  \texttt{doi:10.1002/psp4.12287}. \\
  \textbf{Primary.} Table 2 supplies $\mathrm{CL}/F$, $V/F$ and $K_a$ with RSEs
  and bootstrap CIs; Figure 2 supplies the out-of-sample steady-state check. The
  only primary population-PK source used anywhere in the model.

\item[AndroGel 1\% prescribing information] FDA label, \S12.3 and Table 5. \\
  \textbf{Regulatory primary.} Supplies transdermal $F = 0.10$ verbatim, the
  Day-180 steady-state $C_{\mathrm{avg}}$ by dose arm, and the washout statement
  used for the terminal half-life. AndroGel 1.62\% Table 6 supplies the
  corroborating dose-response inversion discussed in
  Section~\ref{sec:limitations}.

\item[Swerdloff R.S.\ et al.\ (2020) / Jatenzo label] Oral testosterone
  undecanoate. \\
  \textbf{Regulatory primary.} Supplies $C_{\mathrm{avg}}$, $C_{\max}$ and the
  $t_{\max}$ window at 237\,mg BID, with the SD used to argue mean-of-maxima bias.

\item[transfemscience.org compiled tables (9 and 10)] Estradiol ester PK. \\
  \textbf{Secondary, high quality.} A compilation of published single-dose and
  steady-state estradiol data. Used for EV and EC. Not a primary table, and
  flagged as such.

\item[Wang C.\ et al.\ (2010)] Injectable testosterone undecanoate. \\
  \textbf{Secondary.} Cited for the $\sim$1100\,ng/dL single-dose $C_{\max}$
  plausibility check on TU. Not fit to.

\item[Endocrine Society clinical practice guidelines] Referenced for the
  therapeutic reference bands used to score the personas. Not used in any fit.
\end{description}
\endgroup

\paragraph{Source-quality summary.} One primary population-PK paper (TC), two
regulatory labels (T gel, oral T), one high-quality secondary compilation (EV,
EC), and \emph{nothing at all} for TE, TP and TU beyond the clearance identity
that ties them to TC. Table~\ref{tab:tmax} is the honest picture of what that
buys: correct magnitudes everywhere, correct shape for four of the eight forms
with a $t_{\max}$.

\vfill
\noindent\rule{\linewidth}{0.4pt}\\[0.3em]
{\footnotesize This document describes a consumer tracking application's internal
model. It is not clinical guidance, and its predictions are population estimates
that should not be used to make dosing decisions without laboratory
confirmation.}

\end{document}
